\chapter{Preface}

Memory used to be a side effect of how language models compute: the key--value cache of a
Transformer, or the hidden state of a recurrent network. Between 2025 and 2026 it became a
design dimension in its own right. Models now write to memory modules at test time (Titans,
TTT-E2E, In-Place TTT), look facts up in hashed tables (Engram), edit matrix-valued states
with delta rules (Gated DeltaNet, Kimi Delta Attention, Gated DeltaNet-2), filter their memory
like a Kalman filter (Kalman Linear Attention, Gated KalmaNet), and decide per token which
memory to consult (AMOR, HAM). The survey \emph{Memory for Large Language Models}
\citep{zhoubian2026memory} organizes this landscape along three axes: how memory is
represented, when it is updated, and how long it persists.

This book turns that survey into a course. It starts from the standard decoder-only
Transformer and builds, step by step, the mathematics, algorithms and code needed to
understand every 2025--2026 memory system the survey covers. For each system it explains the
idea, derives the equations, gives pseudocode, walks through the authors' implementation, and
provides a short reference implementation written from scratch.

\section*{Who this book is for}

Readers who know linear algebra, basic probability, Python and the fundamentals of deep
learning (multilayer networks, backpropagation, stochastic gradient descent), and who want to
read and modify the code of modern memory architectures. No prior knowledge of state-space
models, linear attention or test-time training is assumed.

\section*{How the book is organized}

\begin{description}
  \item[Part I, Foundations.] \Cref{ch:foundations} sets up the standard LLM: tokenization,
    attention, positional encodings, mixture-of-experts layers, training, the key--value cache
    and the GPU memory hierarchy. \Cref{ch:taxonomy} presents the survey's taxonomy of memory
    and how memory is evaluated. \Cref{ch:linear} develops linear attention, the delta rule,
    state-space models and chunkwise computation, the toolkit for Part~II.
  \item[Part II, Implicit memory.] Memory that lives in the forward computation: sparse and
    selective attention (\cref{ch:sparse}), the delta-rule family (\cref{ch:delta}), more
    expressive state transitions (\cref{ch:expressive}), structured recurrent memory
    (\cref{ch:structured}) and probabilistic or objective-shaped memory
    (\cref{ch:probabilistic}).
  \item[Part III, Explicit memory.] Memory with its own write rule: test-time training
    (\cref{ch:ttt}), memory slots and lookup tables (\cref{ch:lookup}), and multi-timescale
    memory with stable parametric updates (\cref{ch:multiscale}).
  \item[Part IV, Hybrids and memory systems.] Architectures that combine memories
    (\cref{ch:hybrid}), route between them adaptively, and manage the working memory
    (\cref{ch:routing}).
  \item[Part V, Synthesis.] A cross-cutting comparison and the open problems
    (\cref{ch:synthesis}). \Cref{app:code} is a guide to the mirrored code.
\end{description}
Readers new to the area should read Part~I in order. After that, Parts~II--IV can be read in
any order; each chapter points back to the background it relies on.

\section*{Conventions}

Notation is shared across chapters: sequences have length $T$, recurrent memories are
matrices $\mS_t \in \R^{d_v \times d_k}$ written by outer products $\vv_t\vk_t\T$ and read by
$\vo_t = \mS_t\vq_t$, and gates are written $\alpha_t$ (decay) and $\beta_t$ (write strength).
Coloured boxes mark recurring material:
\begin{keyidea}[title={Key idea}]
The central insight of a section.
\end{keyidea}
\begin{taxonomy}
Where a system sits on the survey's three axes and which update rule it uses.
\end{taxonomy}
\begin{implnote}
How the mirrored implementation realizes the equations: files, functions, tensor shapes.
\end{implnote}
\begin{results}
Quantitative results, quoted only from files in the \texttt{mem-sys} repository, with the
source file named.
\end{results}
\begin{pitfall}
A common misunderstanding or implementation trap.
\end{pitfall}
Code listings captioned ``Reference implementation (ours)'' were written for this book to
expose the core mechanism. They favour clarity over speed and are not the authors' code.
Listings captioned ``Excerpt from \ldots'' are quoted from the mirrored repositories.

\section*{About the code and the results}

All implementations discussed are mirrored, at pinned commits, in the \texttt{repos/}
directory of \url{https://github.com/utkrshkmr/mem-sys}; \cref{app:code} lists them. Some
systems have no public code, and some have only unofficial implementations; the text says so
each time. To avoid misquoting, quantitative results are reproduced only when they appear in
a file in the repository (a README, a results table, or a paper PDF included in a mirror).
Otherwise the book describes the authors' claims qualitatively and refers to the paper.

\section*{How this book was written}

This book was drafted with AI assistance (Claude Code) from the survey, the mirrored source
code and the repositories' documentation. Every chapter was compiled and its from-scratch
reference implementations were checked numerically where feasible, but the text has not been
reviewed by the authors of the papers it describes. Treat it as a guide to the primary
sources, and check details against the papers and code before relying on them. Corrections
are welcome as issues or pull requests on the repository.

\section*{Acknowledgements}

This book would not exist without the survey by Sining Zhoubian, Dan Zhang, Evgeny Kharlamov
and Jie Tang, and without the many authors who released their code openly. All credit for the
methods belongs to them; any errors in their exposition are ours.
